\Needspace{11\baselineskip}
\section{Developing the adversary}\label{sec:development}

\subsection{Keeping the comparison graph small}

Start with the larger-half rule. Give each key $x$ a position
$X_x\in(0,1)$ and record a comparison answered $x<y$ by imposing
$X_x<X_y$. Let $P$ be their partial order. Its feasible region is
\[
K(P)=\{X\in(0,1)^n:X_x<X_y\text{ whenever }x<_P y\}.
\]
Each total order occupies volume $1/n!$, so $\vol(K(P))=e(P)/n!$
\cite{stanley}. To decide which answer retains more orders, compare the
volumes on the two sides of $X_x=X_y$.
Dyer, Frieze, and Kannan~\cite{dfk} approximate such volumes in
randomized polynomial time. We need a deterministic way to choose an
answer that keeps almost half the volume.

We record comparisons as arcs and use finite calculations around the
queried vertices. A vertex with many neighbors makes even one
coordinate update depend on many other positions. A long directed path
causes a different problem. In a pass that updates predecessors before
successors, a change can travel through the whole path in one pass.
Bounding both degree and directed path length limits how far a
calculation can reach and how many vertices it needs to visit.
As comparisons accumulate, we need a way to remove vertices while
preserving the answers they record.

\subsection{Fixing a position and continuing in a child}

Choose a point $c$ inside the range of feasible positions of a key $x$.
Fixing $X_x=c$ turns each incoming comparison $y<x$ into the bound
$X_y<c$ and each outgoing comparison $x<z$ into $c<X_z$. We can
remove the vertex of $x$ and keep those bounds.

The key still needs a place in the order. Put it in a block anchored
at $c$. Positions below $c$ precede the whole block, and positions above
$c$ follow it. Several keys can enter the block as long as their
placement respects earlier answers. Give each a fresh position in
$(0,1)$ to describe its order inside the block. That block is a child
node with its own comparison graph, so the same operation can be used
there. Each position describes an item's order at the node that stores
it.

For example, start with $0<X_A<X_B<1$, a triangle of area $1/2$.
Fixing $X_A=1/4$ leaves $X_B\in(1/4,1)$. Give $A$ a fresh child
coordinate $U_A\in(0,1)$. The new feasible assignments form a rectangle
of area $3/4$ (Figure~\ref{fig:section-volume}). The volume grows by
a factor of $3/2$, and both regions represent just the order $A<B$.

\begin{figure}[!htb]
\centering
\begin{tikzpicture}[x=3.25cm,y=3.25cm,>=Latex,font=\small]
\begin{scope}
\node at (.5,1.38) {Before refinement};
\fill[pale] (0,0)--(0,1)--(1,1)--cycle;
\draw[ink!55] (0,0)--(1,1)--(0,1);
\draw[->] (0,0)--(1.12,0) node[right] {$X_A$};
\draw[->] (0,0)--(0,1.12) node[above] {$X_B$};
\draw[dashed,ink!60] (.25,0)--(.25,.25);
\draw[accent,very thick] (.25,.25)--(.25,1);
\node[below] at (.25,0) {$\frac14$};
\node[below left] at (0,0) {$0$};
\node[below] at (1,0) {$1$};
\node[left] at (0,1) {$1$};
\node at (.55,.83) {area $\frac12$};
\end{scope}
\draw[->,thick] (1.2,.53)--(1.63,.53);
\node[above] at (1.415,.56) {refine $A$};
\begin{scope}[shift={(1.9,0)}]
\node at (.5,1.38) {After refinement};
\fill[pale] (0,.25) rectangle (1,1);
\draw[ink!55] (0,.25)--(1,.25)--(1,1)--(0,1);
\draw[->] (0,0)--(1.12,0) node[right] {$U_A$};
\draw[->] (0,0)--(0,1.12) node[above] {$X_B$};
\node[below left] at (0,0) {$0$};
\node[left] at (0,.25) {$\frac14$};
\node[below] at (1,0) {$1$};
\node[left] at (0,1) {$1$};
\node at (.5,.65) {area $\frac34$};
\end{scope}
\end{tikzpicture}
\caption{The highlighted section of the triangle has length $3/4$.
Replacing the fixed parent coordinate by a fresh unit coordinate in the
child gives the rectangle on the right. Its area is $3/2$ times the
original area.}
\label{fig:section-volume}
\end{figure}

The factor $3/2$ is the length of the section at $X_A=1/4$ divided by
the area of the triangle. This ratio is the density of $X_A$ at the
chosen point under the uniform distribution on feasible positions.
The same identity holds for a general move
(Lemma~\ref{lem:section-volume}). It gives us a reason to look for
anchors of high density.

Once keys move into blocks, the initial identity between volume and
order count no longer applies. The counting proof must account for
the locations of the blocks as well as the volume inside each node.
A finer binary location takes more bits to specify. We choose anchors
so that their volume gains help pay for this cost, and bound how
often the adversary needs to move a key.
